% Generated draft: requires mathematical and source review.
\documentclass[11pt]{article}
\usepackage[margin=25mm]{geometry}
\usepackage{fontspec}
\setmainfont{Latin Modern Roman}
\usepackage{amsmath,amssymb,mathtools,unicode-math}
\setmathfont{Latin Modern Math}
\usepackage{enumitem,needspace,xurl,hyperref}
\setlength{\parindent}{0pt}
\setlength{\parskip}{5pt}
\setlength{\emergencystretch}{3em}
\newcommand{\N}{\mathbb{N}}
\newcommand{\Z}{\mathbb{Z}}
\newcommand{\Q}{\mathbb{Q}}
\newcommand{\R}{\mathbb{R}}
\newcommand{\C}{\mathbb{C}}
\newcommand{\F}{\mathbb{F}}
\newcommand{\A}{\mathbb{A}}
\newcommand{\PP}{\mathbb{P}}
\newcommand{\E}{\mathbb{E}}
\newcommand{\eps}{\varepsilon}
\newcommand{\dd}{\,\mathrm{d}}
\begin{document}
\section*{UnsolvedMath ID: 68}
\textbf{Source problem code:} catalogue entry
\subsection{Definitions and mathematical statement}
The supplied problem is DARPA's Mathematical Challenge Four (2007), entitled ``21st Century Fluids.'' Its requested objective is to extend classical fluid dynamics to the following complex-material classes:
\[
\{\text{foams},\ \text{suspensions},\ \text{gels},\ \text{liquid crystals}\}.
\]
The background identifies microstructure, non-Newtonian constitutive behavior, and multiphase dynamics as phenomena that motivate this extension. Here a constitutive law means a mathematical relation specifying a material's response, for example its stress in relation to its deformation and internal state. The source supplies no such law and does not identify a particular fluid model.

\textbf{Source-faithful research objective.} Develop mathematical models and analysis extending classical fluid dynamics to account for complex substances of the listed kinds and their relevant microstructure, non-Newtonian response, or multiphase behavior. The list identifies intended applications; it does not assert that every listed substance exhibits every listed mechanism.

This objective is not yet a quantified theorem or conjecture. To state a concrete mathematical problem, a constitutive law and target material, the unknown fields and governing equations, the physical regime and spatial domain, initial and boundary conditions, admissible data and solution classes, and the desired well-posedness or predictive conclusion must be specified. None of these are selected by the supplied record. In particular it supplies neither a universal equation for all four classes nor a demand for a particular existence, uniqueness, regularity, or convergence theorem.

The literature triage dated 2026-08-17 describes progress in models and analysis for many complex-fluid classes, but classifies the challenge as a research agenda spanning non-equivalent equations and materials. It explicitly says that there is no unitary solved/open status and that a full solution or refutation is not applicable before a target model or theorem is fixed. The registry's \emph{open} label and the triage's \emph{partially solved} label are retained as supplied classifications, not as an assertion about any particular fluid equation.
\subsection{Short English statement}
Develop mathematics for foams, suspensions, gels, and liquid crystals, taking account of phenomena such as microstructure, non-Newtonian behavior, and multiple phases. DARPA's challenge is a broad research objective; the supplied record does not specify the material law, equations, regime, or theorem needed for a single precise problem.
\subsection{Sources}
\begin{itemize}
\item UnsolvedMath local registry, data/unsolvedmath/problems.json, permanent ID 68, source code catalogue entry, ``21st Century Fluids.'' The record attributes the challenge to DARPA (2007).
\item DARPA DSO, Broad Agency Announcement 07-68, Mathematical Challenge Four (2007), supplied reference: \url{https://web.math.utk.edu/~vasili/refs/darpa07.MathChallenges.html}. The embedded review identifies this as the primary challenge text and cites its list of foams, suspensions, gels, and liquid crystals. This draft relies on the supplied record and does not claim independent consultation of the linked text.
\item Ulam AI and the UnsolvedMath Contributors, \emph{UnsolvedMath: A Curated Collection of Open Mathematics Problems}, including the local record's literature triage dated 2026-08-17. Dataset curation and original metadata are supplied under CC BY 4.0; referenced source material retains its own terms.
\end{itemize}
\end{document}
